% Define el tipo de documento: artículo en este caso.
\documentclass[12pt, a4paper, oneside, titlepage]{article}

%----------- Paquetes necesarios para esta memoria -------------------------------------
\usepackage[utf8]{inputenc}
\usepackage[spanish,es-tabla,es-nodecimaldot]{babel}
\usepackage{amsmath}
\usepackage{amsfonts}
\usepackage{amssymb}
\usepackage{csquotes}	% Manejo de simbolos de cita multiples
\usepackage{graphicx}	% Figuras
% Configuracion de interlineado (\onehalfspace y \singlespacing en capitulos)
%\usepackage{setspace}
\usepackage{float}	% Permite que las figuras floten en la posición en la que aparecen en el código
\usepackage{caption}	% Permite dar estilo a los subtitulos de las figuras (captions)
\usepackage{subcaption}	% Permite sub-captions en sub-figuras (figuras dentro de figuras)
\usepackage{wrapfig}
\usepackage[table]{xcolor}
\usepackage{multirow}	% Permite múltiples filas en tablas
\usepackage{emptypage}	% Elimina numeracion de paginas en blanco
\usepackage{enumitem}	% Permite continuar con la numeración entre listas distintas
\usepackage{algorithm}
\usepackage{algorithmic}

%Paquete para gestionar enlaces
\usepackage{hyperref}
\graphicspath{{./capitulos/img/}} %ruta a la carpeta de imaganes, nos evita poner siempre la ruta

%Paquete para gestionar la bibliografía
\usepackage[backend=bibtex, style=numeric, sorting=none, maxbibnames=6]{biblatex}
%---------------------------------------------------------------------------------------

%Datos del alumno y asignatura
\author{Antonio Ariza García}
\title{Introducción a los modelos computacionales: Práctica 1}

%Fecha de creación de la memoria
\date{Curso académico 2020-2021\\Córdoba, \today}

% Cambia el título del índice de algoritmos
%\renewcommand{\listalgorithmname}{Índice de algoritmos}

% Comienzo del documento
\begin{document}

% Cargar la portada
\begin{titlepage}
\centering
\includegraphics[width=0.18\textwidth]{img/logo_uco_sin_texto.png}
\includegraphics[width=0.30\textwidth]{img/logo_uco_epsc.png}\par\vspace{1cm}
{\scshape\LARGE Universidad de Córdoba\par}
{\scshape\Large Escuela Politécnica Superior de Córdoba\par}
\vspace{1cm}
{\scshape\LARGE Ingeniería Informática\par}
{\scshape\Large Especialidad: Computación\par}
{\scshape\Large Cuarto curso. Primer cuatrimestre\par}
\vspace{1.5cm}
{\scshape\LARGE Introducción a los modelos computacionales.\par}
\vspace{1.5cm}
{\huge\bfseries Práctica 3: Redes neuronales de funciones de base radial.\par}
\vspace{1.2cm}
{\Large\itshape Antonio Ariza García\par}
DNI\par
\href{mailto:correo@uco.es}{i62argaa@uco.es}\par
\vfill

{\large Curso académico 2020-2021\\Córdoba, \today\par}
\end{titlepage}

% Cambia el tipo de numeración de páginas a números romanos
\pagenumbering{roman}
% Muestra el índice de contenido
\tableofcontents
\newpage

% Muestra el índice de figuras y lo incluye en el índice de contenido
%\listoffigures
\addcontentsline{toc}{section}{Índice de figuras}
\newpage

% Muestra el índice de tablas y lo incluye en el índice de contenido
\listoftables
\addcontentsline{toc}{section}{Índice de tablas}
\newpage

% Muestra el índice de algoritmos y lo incluye en el índice de contenido
\listofalgorithms
\addcontentsline{toc}{section}{Índice de algoritmos}
\newpage

% Cambia la numeración de páginas a numeración decimal
\pagenumbering{arabic}

% Cambia el espaciado entre párrafos
\setlength{\parskip}{0.7em}

% Cargarmos las demás partes del documento
\section{Introducción} \label{ch:introduccion}

Este documento está realizado en Latex, forma parte de la asignatura Introducción al aprendizaje automático, de la escuela politécnica de Córdoba.
Es una breve introducción a Weka, software utilizado para aprendizaje automático y la minería de datos, está escrito en Java.
\newline
El siguiente documento tiene como objetivo presentarse como guión y entrega de prácticas, puede utilizarse como pequeña guía a la introducción a Weka y aprendizaje automático.
\section{Implementación Algoritmo RBF}
\subsection{Arquitectura de los modelos}

\begin{itemize}
	\item Una capa de entrada con tantas neuronas como variables tenga la base de datos considerada.
	\item Una capa oculta con un número de neuronas a especificar por el usuario.  Siempre tendremos solo una capa oculta, con todas las neuronas de tipo RBF.
	\item Una capa de salida con tantas neuronas como variables de salida tenga la base de datos considerada.
	\begin{itemize}
		\item Si la base de datos es de regresión: Todas las neuronas serán de tipo lineal.
		\item Si la base de datos es de clasificación, todas las neuronas de la capa de salida serán de tipo Softmax.
	\end{itemize}
\end{itemize}

\subsection{Algoritmo Entrenamiento RBF}
\begin{algorithm}[H]
\caption{Entrenamiento RBF Offline}
\begin{algorithmic}[1]
	\STATE centroidesIniciales $\leftarrow$ aleatoriamente $n_1$ patrones 		     (regresion) o $n_1$ patrones de forma estratificada (clasificacion)
	\STATE centroides $\leftarrow$ K-medias (X,$n_1$,centroidesIniciales)
	\STATE $\sigma_j \leftarrow$ (medias de las distancias al resto de centroides)/2
	\STATE Construir la matriz $R_{(Nx(n1+1))}$, donde $R_{ij} = out_j^1(x_i)$ para $j\neq (n_1 +1)$, y $R_{ij} = 1$ para $j =(n_1 +1)$
	\IF {clasificacion}
		\STATE PesosSalida $\leftarrow$ AplicarRegresionLogistica(R,eta)
	\ELSE
		\STATE PesosSalida $\leftarrow$ CalcularPseudoInversa(R)
	\ENDIF
\end{algorithmic}
\end{algorithm}

\begin{itemize}
	\item[1.] Si es clasificación sacamos centroides de forma estratificada, y si es regresión escogemos aleatoriamente $n_1$ patrones
	\item[2.] Ajuste de centros mediante algoritmo de clustering más popular, el K-means, con tantos clusters como neuronas de la RBF.
	\item[3.] Para ajustar los radios de la RBF tomaremos la mitad de la distancia media al resto de centroides. $\sigma_j$ = $\dfrac{1}{2(n_1-1)} * \sum_{i\neq j} \lvert\lvert c_j-c_i \rvert\rvert$
	\item[4.] Construimos la matriz R, donde tenemos el valor de activación de cada neurona para cada patrón.
	\item[5.] Si es Clasificación
	\begin{itemize}
		\item[6.] Aplicamos regresión logística, utilizamos una función a la que le pasamos la matriz R. (puede incluir regularización, parámetro $\eta$).
	\end{itemize}
	\item[7.] Si es Regresión
	\begin{itemize}
		\item[8.] Una vez construida la matriz R, aplicamos la pseudo-inversa de Moore Penrose.
	\end{itemize}
\end{itemize}




\section{Resultados Prácticos}
Probaremos distintas configuraciones de la red neuronal y ejecutaremos cada configuración con cinco semillas (1, 2, 3, 4 y 5). A partir de los resultados obtenidos, se obtendrá la media y la desviación típica del error, el entrenamiento va a guiarse utilizando la función de
entropía cruzada o el MSE.

Para valorar como funciona el algoritmo implementado en esta práctica, emplearemos un total de tres bases de datos:
\begin{itemize}
\item Problema XOR: esta base de datos representa el problema de clasificación no lineal del XOR. Se utilizará el mismo fichero para train y para test.

\item Base de datos divorce: contiene 127 patrones de entrenamiento y 43 patrones de test. La base de datos contiene la respuesta a una serie de preguntas de un conjunto de encuestas en las que se pretende predecir si se va a producir un divorcio en la pareja. Las respuestas a las preguntas de la encuesta se proporcionan en escala de Likert con valores de 0 a 4. Todas
las variables de entrada se consideran numéricas. La base de datos tiene un total de 54 preguntas (por lo tanto, 54 variables de entrada) y dos
categorías (0 no hay divorcio, 1 hay divorcio).

\item Base de datos noMNIST: la base de datos que se utilizará está compuesta por 900 patrones de entrenamiento y 300 patrones de test. Está formada por un conjunto de letras (de la a a la f) escritas con diferentes tipografías o
simbologías.
\end{itemize}

Crearemos una tabla para cada base de datos, en la que se compare la media y desviación típica de cuatro medidas.
\begin{itemize}
\item Error entrenamiento y de test: Se utilizará la función que el usuario haya elegido para  ajustar los pesos (MSE o entropía cruzada)..
\item CCR de entrenamiento y de test.
\end{itemize}

\subsection{Descripción de los distintos parámetros}
A continuación se explicarán cada parámetro que podrá recibir el programa:
\begin{itemize}
\item Argumento t: Indica el nombre del fichero que contiene los datos de entrenamiento a utilizar. Sin este argumento, el programa no puede funcionar.
\item Argumento T: Indica el nombre del fichero que contiene los datos de test a utilizar. Si no se especifica este argumento, utilizar los datos de entrenamiento como test.
\item Argumento i: Indica el número de iteraciones del bucle externo a  realizar. Si no se especifica, utilizar 1000 iteraciones.
\item Argumento l: Indica el número de capas ocultas del modelo de red neuronal. Si no se especifica, utilizar 1 capa oculta.
\item Argumento h: Indica el número de neuronas a introducir en cada una de las capas ocultas. Si no se especifica, utilizar 4 neuronas.
\item Argumento e: Indica el valor del parámetro eta ($\eta$). Por defecto, utilizar $\eta$ = 0,7.
\item Argumento m: Indica el valor del parámetro mu ($\mu$). Por defecto, utilizar $\mu$ = 1.
\item Argumento v: Indica el ratio de patrones de entrenamiento a utilizar como patrones de validación. Por defecto, utilizar v = 0;0.
\item Argumento d: Indica el valor del factor de decremento a utilizar por cada una de las capas. Por defecto, utilizar F = 1.
\/item Argumento o: Booleano que indica si se va a utilizar la versión online. Si no se especifica utilizaremos la versión offline.
\item Argumento f: Indica la función de error que se va a utilizar durante el aprendizaje (0 para el error MSE y 1 para la entropía cruzada). Por defecto, utilizar el error MSE).
\item Argumento s: Booleano que indica si vamos utilizar la función softmax en la capa de salida. Si no se especifica, utilizaremos la función sigmoide.
\item Argumento w: Indica el nombre del fichero en el que se almacenarán la configuración y el valor de los pesos del modelo entrenado.
\end{itemize}


\subsection{Utilizando función error entropía cruzada y la función softmax en la capa de salida (modo offline)}
En esta prueba no emplearemos factor de decremento (F=1) ni conjunto de validación (v= 0;0). En la práctica anterior la mejor arquitectura fue con dos capas ocultas y 100 neuronas en cada una.

\begin{table}[H]
\centering
\vspace{1ex}
\small
\resizebox{\textwidth}{!}{
\begin{tabular}{|c|c|c|c|c|c}
\hline
\textbf{Arquitecturas} & 
\textbf{Train error} &
\textbf{$\sigma$(Train error)} & 
\textbf{Test error} &
\textbf{$\sigma$(Test error} \\ \hline
\rowcolor{green}
\{n:100:100:k\} & 0.001701 & 0.0001286 & 0.001701 & 0.0001286 \\
\hline
\end{tabular}
}
\caption{Errores Base de datos XOR en distintas arquitecturas}
\label{VariacionesDivorce}
\end{table}

Para la BD Seno la mejor arquitectura sería con dos capas ocultas de 64 neuronas cada una, porque aunque la formada por dos capas ocultas y 100 neuronas tenga un error muy parecido, no es así para la desviación típica, pues presenta 5 veces más por lo que puede dar errores mayores.
\begin{table}[H]
\centering
\vspace{1ex}
\small
\resizebox{\textwidth}{!}{
\begin{tabular}{|c|c|c|c|c|c}
\hline
\textbf{Arquitecturas} & 
\textbf{Train error} &
\textbf{$\sigma$(Train error)} & 
\textbf{Test error} &
\textbf{$\sigma$(Test error)} \\ \hline
\{n:2:k\}       & 0.0297282 & 3.08475e-05 & 0.0365526 & 0.000360624 \\
\{n:4:k\}       & 0.029529  & 0.000411965 & 0.0363783 & 0.000223319 \\
\{n:8:k\}       & 0.0294088 & 0.000534317 & 0.03617   & 0.000366612 \\
\{n:32:k\}      & 0.0290472 & 0.000360613 & 0.0360009 & 0.000806839 \\
\{n:64:k\}      & 0.0280052 & 0.000322043 & 0.0352727 & 0.000584671 \\
\{n:100:k\}     & 0.0285983 & 0.000986158 & 0.036288  & 0.000428383 \\
\{n:2:2:k\}     & 0.0297237 & 4.24326e-05 & 0.036248  & 0.000170634 \\
\{n:4:4:k\}     & 0.0297821 & 3.06686e-05 & 0.0362179 & 0.000150539 \\
\{n:8:8:k\}     & 0.0299407 & 0.000118453 & 0.036485  & 0.000630194 \\
\{n:32:32:k\}   & 0.0296178 & 0.00064876  & 0.0361732 & 0.000585751 \\
\rowcolor{green}
\{n:64:64:k\}   & 0.0269406 & 0.00080941  & 0.034039  & 0.000834419 \\
\{n:100:100:k\} & 0.0259549 & 0.00480641  & 0.0340658 & 0.004816921 \\
\hline
\end{tabular}
}
\caption{Errores Base de datos Seno en distintas arquitecturas}
\label{VariacionesDivorce}
\end{table}

En el caso de Quake tenemos 32 neuronas y con una sola capa oculta.
\begin{table}[H]
\centering
\vspace{1ex}
\small
\resizebox{\textwidth}{!}{
\begin{tabular}{|c|c|c|c|c|c}
\hline
\textbf{Arquitecturas} & 
\textbf{Train error} &
\textbf{$\sigma$(Train error)} & 
\textbf{Test error} &
\textbf{$\sigma$(Test error)} \\ \hline
\{n:2:k\}       & 0.0301022 & 7.06987e-05 & 0.027198  & 5.85502e-05 \\
\{n:4:k\}       & 0.029967  & 5.33303e-05 & 0.0270686 & 8.04283e-05 \\
\{n:8:k\}       & 0.029916  & 9.88792e-05 & 0.027052  & 7.48808e-05 \\
\rowcolor{green}
\{n:32:k\}      & 0.0297653 & 2.91166e-05 & 0.026971  & 2.59938e-05 \\
\{n:64:k\}      & 0.0297662 & 3.47861e-05 & 0.0270388 & 3.04479e-05 \\
\{n:100:k\}     & 0.0297446 & 5.52488e-05 & 0.0270129 & 3.55988e-05 \\
\{n:2:2:k\}     & 0.030148  & 4.05299e-05 & 0.0272508 & 4.38807e-05 \\
\{n:4:4:k\}     & 0.0301119 & 3.07848e-05 & 0.0272228 & 3.27004e-05 \\
\{n:8:8:k\}     & 0.0300999 & 7.20564e-05 & 0.0271879 & 8.2355e-05  \\
\{n:32:32:k\}   & 0.0298273 & 8.00825e-05 & 0.0269784 & 4.39127e-05 \\
\{n:64:64:k\}   & 0.0295095 & 7.63962e-05 & 0.0270458 & 5.78271e-05 \\
\{n:100:100:k\} & 0.0294047 & 1.42461e-05 & 0.0271243 & 1.80462e-05 \\
\hline
\end{tabular}
}
\caption{Errores Base de datos Quake en distintas arquitecturas}
\label{VariacionesDivorce}
\end{table}

Para BD Parkinsons la mejor arquitectura está compuesta por dos capas ocultas y 64 neuronas en cada una.
\begin{table}[H]
\centering
\vspace{1ex}
\small
\resizebox{\textwidth}{!}{
\begin{tabular}{|c|c|c|c|c|c}
\hline
\textbf{Arquitecturas} & 
\textbf{Train error} &
\textbf{$\sigma$(Train error)} & 
\textbf{Test error)} &
\textbf{$\sigma$(Test error} \\ \hline
\{n:2:k\}       & 0.0343992 & 9.38514e-06 & 0.0371213 & 9.13403e-05 \\
\{n:4:k\}       & 0.0257571 & 0.000242337 & 0.0255326 & 8.04128e-05 \\
\{n:8:k\}       & 0.0207401 & 0.00149025  & 0.0216091 & 0.0015686 \\
\{n:32:k\}      & 0.0154767 & 0.000518124 & 0.0173277 & 0.000980353 \\
\{n:64:k\}      & 0.0147414 & 0.000632241 & 0.0164462 & 0.000540946 \\
\{n:100:k\}     & 0.0133494 & 0.000380026 & 0.0155019 & 0.000628582 \\
\{n:2:2:k\}     & 0.0315812 & 0.000894328 & 0.032942  & 0.000926471 \\
\{n:4:4:k\}     & 0.0234318 & 0.0014714   & 0.0239546 & 0.00123368 \\
\{n:8:8:k\}     & 0.0146101 & 0.00204368  & 0.0154939 & 0.00246286  \\
\{n:32:32:k\}   & 0.0090358 & 0.00131826  & 0.0112496 & 0.00153763 \\
\rowcolor{green}
\{n:64:64:k\}   & 0.00777673 & 0.00121271 & 0.0101671 & 0.00159433 \\
\{n:100:100:k\} & 0.00878881 & 0.00585147 & 0.010591  & 0.00506543 \\
\hline
\end{tabular}
}
\caption{Errores Base de datos Parkinsons en distintas arquitecturas}
\label{VariacionesDivorce}
\end{table}

\section{Pruebas modificando el Factor de decremento}
Una vez elegida la mejor arquitectura para cada una de las bases de datos, probaremos modificando el factor de decremento con valores \{1,2\}.

Para la mejor arquitectura de la base de datos XOR.
\begin{table}[H]
\centering
\vspace{1ex}
\small
\resizebox{\textwidth}{!}{
\begin{tabular}{|c|c|c|c|c|c|c}
\hline
\textbf{Arquitecturas} & 
\textbf{F} &
\textbf{Train error} &
\textbf{$\sigma$(Train error)} & 
\textbf{Test error} &
\textbf{$\sigma$(Test error)} \\ \hline
\{n:100:100:k\} & 1 & 0.0017019 & 0.00012862 & 0.0017019 & 0.00012862\\
\{n:100:100:k\} & 2 & 0.0145757 & 0.00814214 & 0.0145757 & 0.00814214\\
\hline
\end{tabular}
}
\caption{Variaciones Factor decremento para base datos XOR }
\label{VariacionesDivorce}
\end{table} 

Para la mejor arquitectura de la base de datos Seno.
\begin{table}[H]
\centering
\vspace{1ex}
\small
\resizebox{\textwidth}{!}{
\begin{tabular}{|c|c|c|c|c|c|c}
\hline
\textbf{Arquitecturas} & 
\textbf{F} &
\textbf{Train error} &
\textbf{$\sigma$(Train error)} & 
\textbf{Test error} &
\textbf{$\sigma$(Test error)} \\ \hline
\{n:64:64:k\} & 1 & 0.0269406 & 0.00080941 & 0.0340395  & 0.00083442\\
\{n:64:64:k\} & 2 & 0.0284589 & 0.00039619 & 0.0360283  & 0.00110174\\
\hline
\end{tabular}
}
\caption{Variaciones Factor decremento para base datos Seno }
\label{VariacionesDivorce}
\end{table} 

Para la mejor arquitectura de la base de datos Quake.
\begin{table}[H]
\centering
\vspace{1ex}
\small
\resizebox{\textwidth}{!}{
\begin{tabular}{|c|c|c|c|c|c|c}
\hline
\textbf{Arquitecturas} & 
\textbf{F} &
\textbf{Train error} &
\textbf{$\sigma$(Train error)} & 
\textbf{Test error} &
\textbf{$\sigma$(Test error)} \\ \hline
\{n:32:k\} & 1 & 0.0297653 & 2.91166e-05 & 0.026971  & 2.59938e-05\\
\{n:32:k\} & 2 & 0.0298116 & 2.73964e-05 & 0.0269728 & 2.02602e-05\\
\hline
\end{tabular}
}
\caption{Variaciones Factor decremento para base datos Quake }
\label{VariacionesDivorce}
\end{table} 

Para la mejor arquitectura de la base de datos Parkinson.
\begin{table}[H]
\centering
\vspace{1ex}
\small
\resizebox{\textwidth}{!}{
\begin{tabular}{|c|c|c|c|c|c|c}
\hline
\textbf{Arquitecturas} & 
\textbf{F} &
\textbf{Train error} &
\textbf{$\sigma$(Train error)} & 
\textbf{Test error} &
\textbf{$\sigma$(Test error)} \\ \hline
\{n:64:64:k\} & 1 & 0.0077767  & 0.00121271 & 0.0101671  & 0.00159433\\
\{n:64:64:k\} & 2 & 0.0130969  & 0.0015374  & 0.0130969  & 0.00153748\\
\hline
\end{tabular}
}
\caption{Variaciones Factor decremento para base datos Parkinsons }
\label{VariacionesDivorce}
\end{table} 

Aquí podemos comprobar como el factor de decremento apenas cambia los resultados, pero empeora nuestro modelo en todas las bases de datos.



\end{document}